\documentclass[12pt,a4paper]{report}
\usepackage[utf8]{inputenc}
\usepackage[T1]{fontenc}
\usepackage[french]{babel}
\usepackage{geometry}
\geometry{top=2.5cm,bottom=2.5cm,left=2.5cm,right=2.5cm}
\usepackage{graphicx}
\usepackage{xcolor}
\usepackage{tikz}
\usetikzlibrary{shapes.geometric,arrows.meta,positioning,calc,shadows.blur,fit,backgrounds}
\usepackage{tcolorbox}
\tcbuselibrary{skins,breakable}
\usepackage{fontawesome5}
\usepackage{enumitem}
\usepackage{hyperref}
\hypersetup{colorlinks=true,linkcolor=indigo,urlcolor=indigo,pdfauthor={MonEcole},pdftitle={Manuel --- Application Mobile Parents (PWA)}}
\usepackage{fancyhdr}
\usepackage{tabularx}
\usepackage{booktabs}
\usepackage{float}
\usepackage{caption}

% ====== COULEURS ======
\definecolor{indigo}{HTML}{4338CA}
\definecolor{violet}{HTML}{7C3AED}
\definecolor{amber}{HTML}{D97706}
\definecolor{emerald}{HTML}{059669}
\definecolor{rose}{HTML}{E11D48}
\definecolor{slate}{HTML}{334155}
\definecolor{slatelight}{HTML}{94A3B8}
\definecolor{teal}{HTML}{0D9488}
\definecolor{sky}{HTML}{0EA5E9}
\definecolor{coverbg}{HTML}{1E1B4B}
\definecolor{coveraccent}{HTML}{6366F1}
\definecolor{lightbg}{HTML}{F8FAFC}
\definecolor{warningbg}{HTML}{FEF3C7}
\definecolor{successbg}{HTML}{ECFDF5}
\definecolor{infobg}{HTML}{EFF6FF}
\definecolor{dangerbg}{HTML}{FEF2F2}
\definecolor{purplebg}{HTML}{FAF5FF}
\definecolor{tealbg}{HTML}{F0FDFA}
\definecolor{whatsapp}{HTML}{128C7E}

% ====== TCOLORBOX ======
\newtcolorbox{astuce}[1][]{enhanced,colback=successbg,colframe=emerald,fonttitle=\bfseries\small,title={\faLightbulb\hspace{6pt}Astuce},left=8pt,right=8pt,top=6pt,bottom=6pt,arc=4pt,boxrule=1pt,breakable,#1}
\newtcolorbox{attention}[1][]{enhanced,colback=warningbg,colframe=amber,fonttitle=\bfseries\small,title={\faExclamationTriangle\hspace{6pt}Attention},left=8pt,right=8pt,top=6pt,bottom=6pt,arc=4pt,boxrule=1pt,breakable,#1}
\newtcolorbox{info}[1][]{enhanced,colback=infobg,colframe=indigo,fonttitle=\bfseries\small,title={\faInfoCircle\hspace{6pt}Information},left=8pt,right=8pt,top=6pt,bottom=6pt,arc=4pt,boxrule=1pt,breakable,#1}
\newtcolorbox{danger}[1][]{enhanced,colback=dangerbg,colframe=rose,fonttitle=\bfseries\small,title={\faTimesCircle\hspace{6pt}Important},left=8pt,right=8pt,top=6pt,bottom=6pt,arc=4pt,boxrule=1pt,breakable,#1}
\newtcolorbox{etape}[1][]{enhanced,colback=white,colframe=indigo!60,fonttitle=\bfseries\small,left=10pt,right=10pt,top=8pt,bottom=8pt,arc=6pt,boxrule=1.5pt,breakable,shadow={2pt}{-2pt}{0pt}{black!12},#1}
\newtcolorbox{prerequis}[1][]{enhanced,colback=purplebg,colframe=violet,fonttitle=\bfseries\small,title={\faClipboard\hspace{6pt}Pr\'erequis},left=8pt,right=8pt,top=6pt,bottom=6pt,arc=4pt,boxrule=1pt,breakable,#1}
\newtcolorbox{resultat}[1][]{enhanced,colback=tealbg,colframe=teal,fonttitle=\bfseries\small,title={\faCheck\hspace{6pt}R\'esultat attendu},left=8pt,right=8pt,top=6pt,bottom=6pt,arc=4pt,boxrule=1pt,breakable,#1}

% ====== HEADERS ======
\pagestyle{fancy}\fancyhf{}
\fancyhead[L]{\small\color{slatelight}\leftmark}
\fancyhead[R]{\small\color{indigo}\textbf{MonEcole}}
\fancyfoot[C]{\small\color{slatelight}\thepage}
\renewcommand{\headrulewidth}{0.5pt}
\renewcommand{\footrulewidth}{0.3pt}

% ====== COMMANDES ======
\newcommand{\bouton}[1]{{\colorbox{emerald}{\small\sffamily\bfseries\color{white}\,#1\,}}}
\newcommand{\onglet}[1]{{\colorbox{indigo}{\small\sffamily\bfseries\color{white}\,#1\,}}}
\newcommand{\champ}[1]{\texttt{\color{violet}#1}}
\newcommand{\btnviolet}[1]{{\colorbox{violet}{\small\sffamily\bfseries\color{white}\,#1\,}}}
\newcommand{\btnamber}[1]{{\colorbox{amber}{\small\sffamily\bfseries\color{white}\,#1\,}}}
\newcommand{\btnrose}[1]{{\colorbox{rose}{\small\sffamily\bfseries\color{white}\,#1\,}}}
\newcommand{\fab}[1]{{\colorbox{indigo}{\small\sffamily\bfseries\color{white}\,#1\,}}}

\begin{document}

% ===================================================================
%                         PAGE DE COUVERTURE
% ===================================================================
\begin{titlepage}
\begin{tikzpicture}[remember picture,overlay]
  \fill[coverbg] (current page.south west) rectangle (current page.north east);
  \fill[coveraccent,opacity=0.12] (current page.north east) ++(-4,1) circle (8cm);
  \fill[violet,opacity=0.08] (current page.south west) ++(3,-2) circle (10cm);
  \fill[coveraccent,opacity=0.06] (current page.center) ++(6,-4) circle (5cm);
  \fill[amber,opacity=0.05] (current page.north west) ++(2,-8) circle (4cm);
  \draw[coveraccent,line width=2pt,opacity=0.4] ([yshift=-6cm]current page.north west) ++(2.5,0) -- ++(16,0);
  \node[fill=coveraccent,rounded corners=14pt,minimum size=2cm,inner sep=0pt] at ([yshift=-3.5cm]current page.north) {\color{white}\fontsize{28pt}{28pt}\selectfont\faMobileAlt};
  \node[anchor=north,text width=14cm,align=center] at ([yshift=-6.5cm]current page.north) {
    \color{white}\fontsize{32pt}{36pt}\selectfont\textbf{Manuel d'Utilisation}\\[12pt]
    \fontsize{18pt}{22pt}\selectfont\color{coveraccent!70!white}Application Mobile \textbf{Parents (PWA)}
  };
  \node[anchor=north,text width=14cm,align=center] at ([yshift=-10cm]current page.north) {
    \color{white!60}\fontsize{11pt}{15pt}\selectfont
    Guide complet pour l'application mobile des parents~:\\
    installation, connexion, suivi scolaire, notes,\\
    paiements et messagerie avec les enseignants.
  };
  \draw[coveraccent,line width=1.5pt] ([yshift=-12.5cm]current page.north) ++(-3,0) -- ++(6,0);
  \node[anchor=north,text width=14cm,align=center] at ([yshift=-13.5cm]current page.north) {
    \color{white!50}\small
    \faBuilding\hspace{4pt}Plateforme MonEcole --- Progressive Web App\\[6pt]
    \faCalendar\hspace{4pt}Mai 2026 --- Version 1.0\\[6pt]
    \faBook\hspace{4pt}Documentation Utilisateur --- 7 Sections
  };
  \node[anchor=south,text width=14cm,align=center] at ([yshift=1.5cm]current page.south) {
    \color{white!30}\footnotesize \faCopyright\hspace{3pt} Nexora Tech --- Tous droits r\'eserv\'es
  };
\end{tikzpicture}
\end{titlepage}

\tableofcontents
\newpage

% ===================================================================
\chapter{Pr\'esentation de l'Application}
% ===================================================================

\section{Qu'est-ce que l'application MonEcole Parents~?}

L'application \textbf{MonEcole Parents} est une \textbf{Progressive Web App (PWA)} qui permet aux parents de suivre la scolarit\'e de leurs enfants directement depuis leur t\'el\'ephone, tablette ou ordinateur. Elle fonctionne comme une application native mais sans n\'ecessiter de t\'el\'echargement depuis un Store.

\begin{info}
Une \textbf{PWA} (Progressive Web App) est une application web qui s'installe sur votre appareil comme une application classique. Elle fonctionne en plein \'ecran, peut fonctionner partiellement hors connexion, et se met \`a jour automatiquement.
\end{info}

\section{Fonctionnalit\'es principales}

L'application offre \textbf{6 fonctions cl\'es}~:

\begin{center}
\begin{tikzpicture}[
  card/.style={draw=#1!40, fill=#1!5, rounded corners=8pt, minimum width=4.8cm, minimum height=1.6cm, align=center, font=\scriptsize\bfseries, text=slate, drop shadow={shadow xshift=0.5pt, shadow yshift=-0.5pt, opacity=0.1}},
]
  \node[card=indigo] (c1) at (0,0) {\faIdCard\ \textbf{Informations}\\[2pt]\tiny Fiche d'identit\'e de l'enfant};
  \node[card=sky, right=6pt of c1] (c2) {\faUserEdit\ \textbf{Profil}\\[2pt]\tiny Modifier photo et donn\'ees};
  \node[card=teal, right=6pt of c2] (c3) {\faTachometerAlt\ \textbf{Tableau de bord}\\[2pt]\tiny Pr\'esence, conduite, \'evolution};

  \node[card=violet, below=10pt of c1] (c4) {\faChartBar\ \textbf{Notes}\\[2pt]\tiny \'Evaluations et r\'esultats};
  \node[card=amber, below=10pt of c2] (c5) {\faCreditCard\ \textbf{Paiements}\\[2pt]\tiny Soldes et historique};
  \node[card=rose, below=10pt of c3] (c6) {\faCommentDots\ \textbf{Messages}\\[2pt]\tiny Communication enseignants};
\end{tikzpicture}
\end{center}

\section{Acc\`es \`a l'application}

L'application est accessible via un navigateur web \`a l'adresse~:

\begin{center}
\begin{tikzpicture}
  \fill[indigo!8, rounded corners=8pt] (0,0) rectangle (10,1.2);
  \draw[indigo!30, rounded corners=8pt, line width=1pt] (0,0) rectangle (10,1.2);
  \node[font=\normalsize\bfseries\ttfamily\color{indigo}] at (5,0.6) {https://monecole.pro/parent/};
\end{tikzpicture}
\end{center}

\newpage
% ===================================================================
\chapter{Installer l'Application sur votre Appareil}
% ===================================================================

L'application MonEcole Parents peut \^etre install\'ee sur votre \'ecran d'accueil comme une application native. La proc\'edure varie selon votre appareil.

\section{Sur Android (Google Chrome)}

\begin{etape}[title=Proc\'edure : Installation sur Android]
\begin{enumerate}[leftmargin=1.5cm,label=\textbf{\arabic*.},itemsep=12pt]
  \item \textbf{Ouvrez Google Chrome} et acc\'edez \`a \texttt{https://monecole.pro/parent/}.

  \item \textbf{Attendez la banni\`ere d'installation} --- Au bout de quelques secondes, Chrome affiche en bas de l'\'ecran une banni\`ere \textit{\guillemotleft~Ajouter MonEcole \`a l'\'ecran d'accueil~\guillemotright}. Si elle appara\^it, \textbf{appuyez sur Installer}.

  \item \textbf{Si la banni\`ere n'appara\^it pas}~:
  \begin{itemize}
    \item Appuyez sur les \textbf{trois points verticaux} (\faEllipsisV) en haut \`a droite de Chrome.
    \item S\'electionnez \textbf{\guillemotleft~Ajouter \`a l'\'ecran d'accueil~\guillemotright} (ou \textit{\guillemotleft~Installer l'application~\guillemotright}).
    \item Confirmez en appuyant sur \textbf{Installer}.
  \end{itemize}

  \item \textbf{L'ic\^one MonEcole} appara\^it sur votre \'ecran d'accueil. Vous pouvez d\'esormais lancer l'application en appuyant dessus, comme n'importe quelle autre application.
\end{enumerate}
\end{etape}

\section{Sur iPhone / iPad (Safari)}

\begin{danger}
Sur iOS, l'installation PWA fonctionne \textbf{uniquement avec Safari}. Chrome, Firefox ou d'autres navigateurs sur iOS ne supportent pas l'installation sur l'\'ecran d'accueil.
\end{danger}

\begin{etape}[title=Proc\'edure : Installation sur iOS]
\begin{enumerate}[leftmargin=1.5cm,label=\textbf{\arabic*.},itemsep=12pt]
  \item \textbf{Ouvrez Safari} (le navigateur par d\'efaut d'Apple) et acc\'edez \`a \texttt{https://monecole.pro/parent/}.

  \item \textbf{Appuyez sur le bouton Partager} --- C'est l'ic\^one carr\'e avec une fl\`eche vers le haut (\faShareSquare) situ\'ee en bas de l'\'ecran (iPhone) ou en haut \`a droite (iPad).

  \item \textbf{Faites d\'efiler} le menu qui appara\^it vers le bas et appuyez sur \textbf{\guillemotleft~Sur l'\'ecran d'accueil~\guillemotright} (ic\^one \faPlusSquare).

  \item \textbf{Confirmez le nom} --- Le nom \textit{MonEcole} est pr\'e-rempli. Appuyez sur \textbf{Ajouter} en haut \`a droite.

  \item \textbf{L'ic\^one MonEcole} appara\^it sur votre \'ecran d'accueil. L'application s'ouvrira en plein \'ecran, sans barre d'adresse Safari.
\end{enumerate}
\end{etape}

\begin{astuce}
Apr\`es l'installation sur iOS, l'application s'ouvre en \textbf{mode plein \'ecran} (standalone) avec la barre d'\'etat bleue fonc\'ee de MonEcole. C'est visuellement identique \`a une application native.
\end{astuce}

\section{Sur Ordinateur (Windows, macOS, Linux)}

\begin{etape}[title=Proc\'edure : Installation sur ordinateur]
\begin{enumerate}[leftmargin=1.5cm,label=\textbf{\arabic*.},itemsep=12pt]
  \item \textbf{Ouvrez Google Chrome} (ou Microsoft Edge) et acc\'edez \`a \texttt{https://monecole.pro/parent/}.

  \item \textbf{Rep\'erez l'ic\^one d'installation} --- Dans la barre d'adresse, \`a droite de l'URL, une petite ic\^one (\faDesktop\ ou \faPlusCircle) appara\^it. Cliquez dessus.

  \item \textbf{Confirmez} en cliquant sur \textbf{Installer} dans la bo\^ite de dialogue.

  \item \textbf{L'application s'ouvre} dans une fen\^etre d\'edi\'ee, sans barre d'adresse. Un raccourci est cr\'e\'e~:
  \begin{itemize}
    \item \textbf{Windows}~: dans le menu D\'emarrer et sur le Bureau.
    \item \textbf{macOS}~: dans le Launchpad.
    \item \textbf{Linux}~: dans le lanceur d'applications.
  \end{itemize}
\end{enumerate}
\end{etape}

\begin{info}
Sur \textbf{tous les appareils}, l'application se met \`a jour automatiquement. Vous n'avez rien \`a t\'el\'echarger manuellement --- il suffit de la rouvrir pour b\'en\'eficier des derni\`eres am\'eliorations.
\end{info}

\newpage
% ===================================================================
\chapter{Connexion et Authentification}
% ===================================================================

\section{Se connecter pour la premi\`ere fois}

L'authentification utilise un syst\`eme de \textbf{code OTP} (One-Time Password) envoy\'e par SMS. Aucun mot de passe \`a retenir.

\begin{prerequis}
\begin{itemize}
  \item Votre \textbf{num\'ero de t\'el\'ephone} doit avoir \'et\'e enregistr\'e par l'\'ecole dans le syst\`eme MonEcole (au moment de l'inscription de votre enfant).
  \item Votre t\'el\'ephone doit pouvoir \textbf{recevoir des SMS}.
\end{itemize}
\end{prerequis}

\begin{etape}[title=Proc\'edure : Se connecter]
\begin{enumerate}[leftmargin=1.5cm,label=\textbf{\arabic*.},itemsep=12pt]
  \item \textbf{Ouvrez l'application} ou acc\'edez \`a \texttt{https://monecole.pro/parent/}.

  \item \textbf{S\'electionnez votre pays} dans le menu d\'eroulant \champ{Pays} (ex.~: Burundi, RDC, Rwanda...). Un drapeau et le pr\'efixe t\'el\'ephonique (+257, +243...) s'affichent automatiquement.

  \item \textbf{S\'electionnez votre \'etablissement} dans le menu d\'eroulant \champ{\'Etablissement}. La liste ne montre que les \'ecoles du pays s\'electionn\'e.

  \item \textbf{Saisissez votre num\'ero de t\'el\'ephone} dans le champ \champ{T\'el\'ephone}. Entrez le num\'ero sans le pr\'efixe pays (ex.~: 79123456).

  \item \textbf{Cliquez sur} \bouton{Envoyer le code}. Un SMS contenant un code \`a 6 chiffres est envoy\'e \`a votre t\'el\'ephone.

  \item \textbf{Saisissez le code OTP} re\c{c}u dans le champ qui appara\^it. Le code est valide pendant \textbf{5 minutes}.

  \item \textbf{Cliquez sur} \bouton{V\'erifier}. Si le code est correct, vous \^etes redirig\'e vers la page d'accueil.
\end{enumerate}
\end{etape}

\begin{attention}
Si vous ne recevez pas le SMS, v\'erifiez que votre num\'ero est correctement enregistr\'e aupr\`es de l'\'ecole. Vous pouvez aussi cliquer sur \textbf{Renvoyer le code} apr\`es 60 secondes.
\end{attention}

\section{Se d\'econnecter}

Pour vous d\'econnecter, appuyez sur le \textbf{bouton flottant} (\faTh) puis s\'electionnez \textbf{D\'econnexion} (\faSignOutAlt\ gris).

\newpage
% ===================================================================
\chapter{Page d'Accueil --- Mes Enfants}
% ===================================================================

\section{Pr\'esentation de l'accueil}

Apr\`es connexion, la page d'accueil affiche~:

\begin{enumerate}[leftmargin=1.2cm]
  \item \textbf{Un message de bienvenue} personnalis\'e avec votre nom.
  \item \textbf{Deux compteurs}~: nombre d'enfants et nombre d'\'etablissements.
  \item \textbf{La liste de vos enfants} sous forme de cartes cliquables. Chaque carte affiche~:
  \begin{itemize}
    \item L'\textbf{avatar} (photo ou initiales avec couleur bleu/rose selon le genre).
    \item Le \textbf{nom et pr\'enom} de l'enfant.
    \item La \textbf{classe} actuelle (ex.~: \textit{3\`eme Ann\'ee Secondaire}).
    \item Le nom de l'\textbf{\'etablissement}.
    \item Le \textbf{pays} (si multi-pays).
  \end{itemize}
\end{enumerate}

\begin{info}
Si vos enfants sont inscrits dans \textbf{diff\'erents \'etablissements} ou m\^eme \textbf{diff\'erents pays}, ils apparaissent tous sur cette m\^eme page. L'architecture multi-tenant de MonEcole g\`ere cette situation automatiquement.
\end{info}

\begin{etape}[title=Proc\'edure : Acc\'eder au dossier d'un enfant]
\begin{enumerate}[leftmargin=1.5cm,label=\textbf{\arabic*.},itemsep=8pt]
  \item \textbf{Rep\'erez la carte} de l'enfant dans la liste.
  \item \textbf{Appuyez sur la carte}. Vous \^etes redirig\'e vers la page de suivi de cet enfant.
  \item \textbf{Pour revenir \`a la liste}, appuyez sur la fl\`eche retour (\faArrowLeft) en haut \`a gauche.
\end{enumerate}
\end{etape}

\newpage
% ===================================================================
\chapter{Le Menu Flottant (FAB)}
% ===================================================================

\section{Fonctionnement}

La navigation dans la fiche de chaque enfant se fait via un \textbf{bouton flottant} (FAB --- Floating Action Button) situ\'e en bas \`a droite de l'\'ecran. C'est un bouton rond indigo avec l'ic\^one \faTh.

\begin{etape}[title=Proc\'edure : Naviguer avec le menu flottant]
\begin{enumerate}[leftmargin=1.5cm,label=\textbf{\arabic*.},itemsep=10pt]
  \item \textbf{Appuyez sur le bouton} rond \fab{\faTh} en bas \`a droite. Un voile semi-transparent appara\^it et le menu s'ouvre vers le haut.

  \item \textbf{Choisissez la section} en appuyant sur le bouton color\'e correspondant. Le menu se referme automatiquement et la section s'affiche.

  \item \textbf{Pour fermer sans naviguer}, appuyez sur le voile ou sur le bouton FAB (devenu rouge avec un \faTimes).
\end{enumerate}
\end{etape}

\section{Les 7 options du menu}

\begin{center}
\begin{tabularx}{\textwidth}{p{1cm}lllX}
\toprule
\textbf{Ic\^one} & \textbf{Label} & \textbf{Couleur} & \textbf{Section} & \textbf{Description} \\
\midrule
\faHome & Accueil & Vert & Page d'accueil & Retour \`a la liste des enfants \\
\faIdCard & Infos & Vert & Informations & Fiche d'identit\'e de l'enfant \\
\faUserEdit & Profil & Bleu & Profil & Modifier les donn\'ees et la photo \\
\faTachometerAlt & Tableau & Teal & Dashboard & Pr\'esence, conduite, \'evolution \\
\faChartBar & Notes & Violet & Notes & \'Evaluations et r\'esultats scolaires \\
\faCreditCard & Paiements & Ambre & Paiements & Soldes et historique des paiements \\
\faCommentDots & Messages & Rose & Messagerie & Communication avec les enseignants \\
\faSignOutAlt & D\'econnexion & Gris & --- & Fermer la session \\
\bottomrule
\end{tabularx}
\end{center}

% manuel_pwa_parents_suite.tex --- inclus par manuel_pwa_parents.tex

\newpage
% ===================================================================
\chapter{Section Informations}
% ===================================================================

\section{Contenu de la fiche}

La section \textbf{Informations} (\faIdCard) affiche les donn\'ees administratives de l'enfant sous forme de grille de cartes~:

\begin{center}
\begin{tabularx}{\textwidth}{lX}
\toprule
\textbf{Champ} & \textbf{Description} \\
\midrule
Nom & Nom de famille de l'\'el\`eve \\
Pr\'enom & Pr\'enom de l'\'el\`eve \\
Genre & Masculin ou F\'eminin \\
Date de naissance & Format jour/mois/ann\'ee \\
Matricule & Num\'ero matricule si attribu\'e par l'\'ecole \\
Code & Code \'el\`eve interne \\
Classe & Classe actuelle (pleine largeur) \\
\bottomrule
\end{tabularx}
\end{center}

\begin{info}
Cette section est en \textbf{lecture seule}. Pour modifier les informations, utilisez la section \textbf{Profil} (bouton \faUserEdit\ dans le menu flottant).
\end{info}

\newpage
% ===================================================================
\chapter{Section Profil --- Modifier les Donn\'ees}
% ===================================================================

\section{Modifier le profil de l'enfant}

\begin{etape}[title=Proc\'edure : Modifier le profil]
\begin{enumerate}[leftmargin=1.5cm,label=\textbf{\arabic*.},itemsep=12pt]
  \item \textbf{Ouvrez le menu flottant} et appuyez sur \textbf{Profil} (\faUserEdit\ bleu).

  \item \textbf{La photo de profil} appara\^it en haut. Appuyez sur \textbf{Changer la photo} pour s\'electionner une image depuis votre galerie ou cam\'era.

  \item \textbf{Modifiez les champs} souhait\'es~:
  \begin{itemize}
    \item \champ{Nom}, \champ{Pr\'enom}, \champ{Post-nom} --- texte libre.
    \item \champ{Genre} --- menu d\'eroulant (Masculin / F\'eminin).
    \item \champ{Date de naissance} --- s\'electeur de date.
    \item \champ{Lieu de naissance} --- texte libre.
    \item \champ{Adresse} --- adresse du domicile.
  \end{itemize}

  \item \textbf{Appuyez sur} \bouton{Enregistrer les modifications} (bouton indigo en bas de page).

  \item \textbf{Un toast vert} confirme la sauvegarde~: \textit{\guillemotleft~Profil mis \`a jour avec succ\`es~\guillemotright}.
\end{enumerate}
\end{etape}

\begin{attention}
Certains champs (comme le \textbf{Matricule} et le \textbf{Code \'el\`eve}) sont g\'er\'es par l'\'ecole et ne sont pas modifiables par le parent. Ils apparaissent gris\'es.
\end{attention}

\newpage
% ===================================================================
\chapter{Section Tableau de Bord}
% ===================================================================

\section{Vue d'ensemble}

Le \textbf{Tableau de bord} (\faTachometerAlt) offre une vue synth\'etique de la vie scolaire de l'enfant avec trois panneaux~:

\subsection{Panneau Pr\'esence}

Trois indicateurs en haut~:

\begin{center}
\begin{tabularx}{\textwidth}{llX}
\toprule
\textbf{Indicateur} & \textbf{Couleur} & \textbf{Description} \\
\midrule
Taux & Neutre & Pourcentage de pr\'esence (ex.~: 94\%) \\
Pr\'esences & Neutre & Nombre de jours pr\'esent \\
Absences & Rouge & Nombre de jours absent \\
\bottomrule
\end{tabularx}
\end{center}

En dessous, la \textbf{liste des absences} d\'etaille chaque jour d'absence avec la date et le motif (justifi\'e ou non).

\subsection{Panneau Conduite}

Affiche les observations de conduite signal\'ees par les enseignants (retards, comportement, remarques positives).

\subsection{Panneau \'Evolution}

Pr\'esente l'\'evolution des r\'esultats scolaires de l'enfant au fil des p\'eriodes.

\newpage
% ===================================================================
\chapter{Section Notes et \'Evaluations}
% ===================================================================

\section{Consulter les notes}

La section \textbf{Notes} (\faChartBar) affiche toutes les \'evaluations de l'enfant et ses r\'esultats. Les donn\'ees sont organis\'ees par \textbf{cours}.

\begin{etape}[title=Proc\'edure : Consulter les notes]
\begin{enumerate}[leftmargin=1.5cm,label=\textbf{\arabic*.},itemsep=10pt]
  \item \textbf{Ouvrez le menu flottant} et appuyez sur \textbf{Notes} (\faChartBar\ violet).

  \item \textbf{La liste des cours} s'affiche. Pour chaque cours, vous voyez~:
  \begin{itemize}
    \item Le \textbf{nom du cours} et l'\textbf{enseignant} responsable.
    \item La liste des \textbf{\'evaluations} avec pour chacune~:
    \begin{itemize}
      \item Le \textbf{titre} de l'\'evaluation (ex.~: \textit{Interrogation Chapitre 3}).
      \item Le \textbf{type} (IE, DS, TP, EX...).
      \item La \textbf{note obtenue} / maximum (ex.~: 15/20).
      \item Un \textbf{badge color\'e} indiquant la performance~:
      \begin{itemize}
        \item \textcolor{emerald}{\textbf{Vert}} --- bonne note ($\geq$ 70\%).
        \item \textcolor{amber}{\textbf{Orange}} --- note moyenne ($\geq$ 50\%).
        \item \textcolor{rose}{\textbf{Rouge}} --- note insuffisante ($<$ 50\%).
      \end{itemize}
    \end{itemize}
  \end{itemize}

  \item \textbf{Les pi\`eces jointes} --- Si l'\'evaluation a un document attach\'e (sujet PDF), un bouton permet de le t\'el\'echarger directement.
\end{enumerate}
\end{etape}

\begin{info}
Les notes affich\'ees sont uniquement les notes \textbf{valid\'ees} par l'\'ecole. Si une note n'appara\^it pas encore, c'est que l'enseignant ne l'a pas encore publi\'ee.
\end{info}

\begin{astuce}
Les \'evaluations sont distingu\'ees par un \textbf{badge de port\'ee}~:
\begin{itemize}
  \item \textcolor{violet}{\textbf{Personnalis\'e}} --- \'evaluation concernant uniquement votre enfant (ou note individuelle).
  \item \textcolor{sky}{\textbf{Collectif}} --- \'evaluation commune \`a toute la classe.
\end{itemize}
\end{astuce}

\newpage
% ===================================================================
\chapter{Section Paiements}
% ===================================================================

\section{Consulter la situation financi\`ere}

La section \textbf{Paiements} (\faCreditCard) affiche deux panneaux~:

\subsection{Panneau Soldes}

Pour chaque variable de paiement (frais scolaires, uniforme, transport...), une carte affiche~:
\begin{itemize}
  \item Le \textbf{nom de la variable} et sa \textbf{cat\'egorie}.
  \item Une \textbf{barre de progression} color\'ee indiquant le pourcentage pay\'e.
  \item Le \textbf{montant pay\'e} et le \textbf{reste \`a payer}.
  \item La \textbf{r\'eduction} \'eventuelle (en pourcentage).
\end{itemize}

\subsection{Panneau Historique}

Liste chronologique de tous les paiements effectu\'es avec~:
\begin{itemize}
  \item La \textbf{date} du paiement.
  \item Le \textbf{montant} en Fbu.
  \item La \textbf{variable} concern\'ee.
  \item Le \textbf{statut}~: \textcolor{emerald}{Valid\'e}, \textcolor{amber}{En attente}, ou \textcolor{rose}{Rejet\'e}.
\end{itemize}

\begin{etape}[title=Proc\'edure : V\'erifier les paiements]
\begin{enumerate}[leftmargin=1.5cm,label=\textbf{\arabic*.},itemsep=8pt]
  \item \textbf{Ouvrez le menu flottant} et appuyez sur \textbf{Paiements} (\faCreditCard\ ambre).
  \item \textbf{Consultez les soldes} en haut --- chaque carte montre la progression du paiement.
  \item \textbf{D\'efilez vers le bas} pour voir l'\textbf{historique} des transactions pass\'ees.
\end{enumerate}
\end{etape}

\newpage
% ===================================================================
\chapter{Section Messagerie}
% ===================================================================

\section{Pr\'esentation de la messagerie}

La section \textbf{Messages} (\faCommentDots) offre une messagerie interne style \textbf{WhatsApp} permettant la communication bidirectionnelle entre le parent et les enseignants / la direction de l'\'ecole.

\begin{info}
Lorsque vous ouvrez la messagerie, l'interface passe en \textbf{mode immersif}~: le profil de l'enfant dispara\^it et la messagerie occupe tout l'\'ecran, exactement comme WhatsApp.
\end{info}

\section{Interface de la messagerie}

L'interface se compose de~:

\begin{enumerate}[leftmargin=1.2cm]
  \item \textbf{Barre verte} en haut --- titre \textit{\guillemotleft~Messages~\guillemotright} avec le nom de l'enfant.
  \item \textbf{Barre de recherche} --- pour filtrer les contacts par nom.
  \item \textbf{Liste des contacts} organis\'ee en deux groupes~:
  \begin{itemize}
    \item \textcolor{amber}{\textbf{Direction}} --- personnel administratif (directeur, pr\'efet d'\'etudes).
    \item \textcolor{violet}{\textbf{Enseignants}} --- enseignants de la classe de l'enfant, avec le nom du cours entre parenth\`eses.
  \end{itemize}
  \item \textbf{Zone de conversation} --- bulles de messages avec horodatage.
  \item \textbf{Barre de saisie} en bas --- champ texte + bouton d'envoi vert.
\end{enumerate}

\section{Pour lire et r\'epondre aux messages}

\begin{etape}[title=Proc\'edure : Lire et r\'epondre]
\begin{enumerate}[leftmargin=1.5cm,label=\textbf{\arabic*.},itemsep=12pt]
  \item \textbf{Ouvrez le menu flottant} et appuyez sur \textbf{Messages} (\faCommentDots\ rose).

  \item \textbf{S\'electionnez un contact} dans la liste \`a gauche. Les contacts avec des messages non lus affichent un \textbf{badge vert} avec le nombre de messages.

  \item \textbf{La conversation s'ouvre} avec les messages pr\'ec\'edents. Les messages envoy\'es par vous sont align\'es \`a \textbf{droite} (fond vert clair). Les messages re\c{c}us sont \`a \textbf{gauche} (fond blanc).

  \item \textbf{Pour r\'epondre}, tapez votre message dans le champ en bas et appuyez sur le \textbf{bouton vert} (\faPaperPlane).

  \item \textbf{Le message appara\^it} imm\'ediatement dans la conversation. L'enseignant re\c{c}oit \'egalement une \textbf{notification par e-mail} pour \^etre averti.
\end{enumerate}
\end{etape}

\section{Types de messages}

Les messages sont distingu\'es par leur \textbf{port\'ee}~:

\begin{center}
\begin{tabularx}{\textwidth}{lX}
\toprule
\textbf{Type} & \textbf{Description} \\
\midrule
\textcolor{violet}{\textbf{Personnalis\'e}} & Message adress\'e sp\'ecifiquement au parent concernant son enfant. \\
\textcolor{sky}{\textbf{Collectif}} & Message envoy\'e par l'enseignant \`a tous les parents de la classe. \\
\bottomrule
\end{tabularx}
\end{center}

\begin{astuce}
Les messages \textbf{collectifs} sont identifiables par un badge bleu \textit{\guillemotleft~Classe enti\`ere~\guillemotright}. Les messages \textbf{personnalis\'es} portent le badge violet \textit{\guillemotleft~Personnalis\'e~\guillemotright}.
\end{astuce}

\begin{attention}
Lorsque vous r\'epondez, votre message est \textbf{toujours individuel} (personnalis\'e). Seul l'enseignant concern\'e le voit, m\^eme si le message initial \'etait collectif.
\end{attention}

\newpage
% ===================================================================
\chapter{R\'ef\'erence Rapide}
% ===================================================================

\section{Raccourcis de navigation}

\begin{tabularx}{\textwidth}{llX}
\toprule
\textbf{Geste} & \textbf{\'El\'ement} & \textbf{Action} \\
\midrule
Appui sur FAB (\faTh) & Bouton flottant & Ouvre/ferme le menu de navigation \\
Appui sur carte enfant & Page d'accueil & Ouvre la fiche de l'enfant \\
\faArrowLeft\ en haut & Header & Retourne \`a la page pr\'ec\'edente \\
\faSignOutAlt\ en haut & Header & D\'econnexion \\
\bottomrule
\end{tabularx}

\section{R\'esolution de probl\`emes courants}

\begin{center}
\begin{tabularx}{\textwidth}{lX}
\toprule
\textbf{Probl\`eme} & \textbf{Solution} \\
\midrule
Je ne re\c{c}ois pas le code OTP & V\'erifiez votre num\'ero aupr\`es de l'\'ecole. Attendez 60s et cliquez \textit{Renvoyer}. \\
Aucun enfant n'appara\^it & Votre num\'ero n'est pas associ\'e \`a un \'el\`eve. Contactez l'administration. \\
Les notes sont vides & L'enseignant n'a pas encore publi\'e les r\'esultats. \\
L'appli ne s'installe pas (iOS) & Utilisez \textbf{Safari} uniquement. Chrome sur iOS ne supporte pas les PWA. \\
Messages envoy\'es disparaissent & Videz le cache du navigateur ou red\'emarrez l'application. \\
\bottomrule
\end{tabularx}
\end{center}

\section{Glossaire}

\begin{description}[leftmargin=1.5cm, style=nextline, font=\bfseries\color{indigo}]
  \item[FAB] Floating Action Button --- le bouton rond flottant en bas \`a droite servant de menu de navigation.
  \item[OTP] One-Time Password --- code \`a usage unique envoy\'e par SMS pour l'authentification.
  \item[PWA] Progressive Web App --- application web installable sur l'\'ecran d'accueil sans passer par un Store.
  \item[Service Worker] Programme en arri\`ere-plan qui g\`ere le cache et les mises \`a jour automatiques de l'application.
  \item[Toast] Notification temporaire affich\'ee en haut de l'\'ecran pour confirmer une action.
\end{description}

\vspace{2cm}
\begin{center}
\begin{tikzpicture}
  \fill[coverbg, rounded corners=12pt] (0,0) rectangle (14,2.5);
  \node[text=white, font=\large\bfseries] at (7,1.5) {\faCheck\hspace{8pt}Fin du Manuel --- Application Parents (PWA)};
  \node[text=white!50, font=\small] at (7,0.7) {Plateforme MonEcole --- Nexora Tech --- Mai 2026 --- v1.0};
\end{tikzpicture}
\end{center}


\end{document}
